\documentclass[11pt,a4paper]{article}
\usepackage[margin=1in]{geometry}
\usepackage[T1]{fontenc}
\usepackage[utf8]{inputenc}
\usepackage{lmodern}
\usepackage{microtype}
\usepackage{graphicx}
\usepackage{amsmath}
\usepackage{hyperref}
\usepackage{parskip}
\usepackage{enumitem}

\hypersetup{
    colorlinks=true,
    linkcolor=blue,
    urlcolor=blue,
    pdftitle={A Reader's Guide to Decoherence as Dimensional Collapse},
    pdfauthor={Ian Todd}
}

\title{\textbf{A Reader's Guide to\\Decoherence as Dimensional Collapse}}
\author{Ian Todd}
\date{Companion note for the manuscript\\\emph{Decoherence as Dimensional Collapse: The Geometry of the Quantum-Classical Transition}}

\begin{document}

\maketitle

\section*{Who this is for}

This note is for readers who want to understand what the paper proves without working through 37 pages of proofs.  It assumes you are comfortable with science writing and know roughly what a quantum state is, but not that you have met information geometry before.  Every technical term is explained when it first appears.

It is not a substitute for the paper.  It is a map: what question the paper asks, what it proves, what is genuinely new, and what remains open.

\section*{The one-sentence claim}

The central claim is that \textbf{decoherence---the process by which quantum systems lose their quantum character---is not just the suppression of interference.  It is a structured geometric projection from a higher-dimensional quantum state space onto a lower-dimensional classical one.}

That statement has three parts:
\begin{enumerate}[leftmargin=1.5em]
\item quantum states have more distinguishable directions than classical probability distributions;
\item dephasing removes exactly the directions associated with quantum coherence;
\item the geometry of the surviving classical directions can be quantified very precisely.
\end{enumerate}

\section*{Background: what you need to know}

\subsection*{Quantum states and density matrices}

A quantum system with $d$ energy levels is described by a \emph{density matrix}~$\rho$: a $d \times d$ matrix that is positive (all eigenvalues non-negative) and has trace~1.  The diagonal entries of~$\rho$ are the \emph{populations}---the probabilities of finding the system in each energy level.  The off-diagonal entries are the \emph{coherences}---they encode phase relationships between energy levels and are responsible for interference effects.

A \emph{classical} probability distribution is just the diagonal: $d$ non-negative numbers summing to~1.  So a classical state is a special case of a quantum state---one with all coherences set to zero.

\subsection*{Decoherence and dephasing}

When a quantum system interacts with its environment (air molecules, photons, vibrations), its coherences decay.  This is \emph{decoherence}.  \emph{Dephasing} is the idealized version: the coherences are set to zero while the populations are left unchanged.

After dephasing, the quantum state becomes a classical probability distribution.  This is the quantum-to-classical transition, and it happens extraordinarily fast for macroscopic objects---Zurek estimates $10^{-31}$ seconds for a dust grain in a photon bath.

\subsection*{Degrees of freedom}

A $d \times d$ density matrix has $d^2 - 1$ independent real parameters (the $-1$ comes from the trace constraint).  Of these, $d - 1$ are populations and $d^2 - d$ are coherences.  For a qubit ($d = 2$): 3 total parameters, 1 population, 2 coherences.  For a qutrit ($d = 3$): 8 total, 2 populations, 6 coherences.

Dephasing eliminates the $d^2 - d$ coherence parameters and preserves the $d - 1$ population parameters.  That is the dimensional collapse:
\[
d^2 - 1 \;\longrightarrow\; d - 1.
\]

\subsection*{Metrics: measuring distances between states}

To make the ``geometry'' precise, we need a way to measure distances between quantum states.  A \emph{metric} assigns a notion of distance (or angle) to nearby states.

The \textbf{Bures metric} is the standard distance measure on quantum states.  Intuitively, it measures how statistically distinguishable two nearby states are: if two states are hard to tell apart by any measurement, their Bures distance is small.  It is the unique ``smallest'' metric with this property (technically: the minimal monotone Riemannian metric under quantum channels).

The \textbf{Fisher--Rao metric} is the analogous distance measure on classical probability distributions.  It measures how distinguishable two nearby probability distributions are.

A key fact, known before this paper: when you restrict the Bures metric to diagonal states (classical probability distributions), you recover the Fisher--Rao metric exactly.  So the classical metric lives naturally inside the quantum metric.

\subsection*{Tangent spaces and orthogonality}

At any state~$\rho$, the \emph{tangent space} is the set of all directions you can move in state space.  It splits naturally into two parts:
\begin{itemize}[leftmargin=1.5em]
\item \emph{population directions}: changes to the diagonal entries;
\item \emph{coherence directions}: changes to the off-diagonal entries.
\end{itemize}

The paper asks: are these two sets of directions \emph{perpendicular} (orthogonal) under the Bures metric?  If they are, then classical and quantum information are geometrically independent---you can change one without affecting the other.  If they are not, the two types of information are entangled in the metric geometry.

\section*{The basic picture}

\begin{figure}[ht]
\centering
\includegraphics[width=0.92\textwidth]{figures/fig1_bloch_collapse.pdf}
\caption{Dephasing removes coherence directions and leaves the classical simplex of populations.  For a qubit, the full quantum state space is the 3-dimensional Bloch ball; dephasing collapses it onto the 1-dimensional $z$-axis (the classical populations).}
\end{figure}

The paper uses the phrase \emph{dimensional collapse} because this is not just noise.  Dephasing does not blur the state randomly---it projects it onto a specific, lower-dimensional subspace.  The dimensions that are lost are precisely the coherence directions, and the dimensions that survive are precisely the population directions.

\section*{What is standard, and what is new?}

The paper is explicit about this.

\subsection*{Mostly standard, but organized into one framework}

\begin{itemize}[leftmargin=1.5em]
\item Dephasing kills off-diagonal matrix elements.
\item The Bures metric is the natural distance measure on quantum states.
\item On diagonal states, the Bures metric reduces to the Fisher--Rao metric.
\item The free-energy identity for dephasing is known once coherence is treated as a thermodynamic resource.
\end{itemize}

Those ingredients were already in the literature, but not assembled into an explicit geometric picture of the quantum-to-classical transition.

\subsection*{The genuinely new parts}

\begin{enumerate}[leftmargin=1.5em]
\item \textbf{Exact qubit angle theorem.}  At every interior qubit state, the paper computes a closed-form formula for the angle between the population direction and the coherence directions.  This angle tells you how cleanly separated the classical and quantum degrees of freedom are.

\item \textbf{Operational corollary.}  That angle is not just geometric decoration.  It has a concrete statistical meaning: it tells you how much \emph{Fisher information} (a measure of how precisely you can estimate a parameter from data) about the population survives when coherence is treated as an unknown nuisance.

\item \textbf{Spectator theorem ($d \geq 3$).}  In three or more dimensions, the paper proves that if two energy levels are coherent but some third level is uninvolved (``decoupled''), then the population and coherence directions are \emph{not} perpendicular.  The uninvolved level acts as a geometric witness---its contribution to the metric stays zero while the coherent levels' contributions become nonzero, breaking the symmetry.

\item \textbf{Perturbative theorem ($d \geq 3$).}  Near the classical probability simplex (i.e., states with small coherence), any nonzero coherence breaks perpendicularity for $d \geq 3$.  The paper gives an explicit formula for how much.
\end{enumerate}

\subsection*{Still open}

The paper does \emph{not} prove the full global statement: that for every interior state in every dimension $d \geq 3$, perpendicularity holds \emph{only} at diagonal states.  That remains a conjecture.  The spectator theorem proves it for states where at least one coherence pair has a decoupled third level; the perturbative theorem proves it near the simplex; and an adversarial numerical search (100 optimization runs per dimension, actively trying to find counterexamples) finds none.  But a complete proof for fully-coherent states remains open.

\section*{The main result: the angle formula}

At diagonal states, the population and coherence directions are perfectly perpendicular ($90^\circ$) under the Bures metric.  This is the cleanest possible separation.

Away from the diagonal, the separation becomes oblique---the angle drops below $90^\circ$.  For qubits, the paper proves exactly:
\[
\cos^2 \theta_B
=
\frac{r_z^2\, r_\perp^2}{(1-r_\perp^2)(1-r_z^2)},
\]
where:
\begin{itemize}[leftmargin=1.5em]
\item $r_z$ is the \emph{population imbalance}---how unequally the two energy levels are occupied (zero means equal populations);
\item $r_\perp = \sqrt{r_x^2 + r_y^2}$ is the \emph{coherence magnitude}---how much quantum interference is present (zero means diagonal);
\item $|\mathbf{r}| < 1$ ensures the state is in the interior of the Bloch ball (not a pure state on the surface).
\end{itemize}

This formula says:
\begin{itemize}[leftmargin=1.5em]
\item the split is perfectly perpendicular ($\theta_B = 90^\circ$) when the state is diagonal ($r_\perp = 0$);
\item it is also perpendicular when populations are equal ($r_z = 0$), even if coherence is present;
\item it becomes oblique only when the state has \emph{both} nonzero coherence \emph{and} unequal populations;
\item under dephasing ($r_\perp$ shrinking toward zero), the angle moves back toward $90^\circ$.
\end{itemize}

This is the key conceptual point: \textbf{the metric separation between classical and quantum information is state-dependent, and dephasing sharpens it.}  As decoherence proceeds, the classical/quantum split becomes more and more perpendicular---until at full dephasing the two kinds of information are perfectly separated.

\begin{figure}[ht]
\centering
\includegraphics[width=0.92\textwidth]{figures/fig5_bures_angle.pdf}
\caption{The qubit Bures angle across the Bloch ball.  Orthogonality ($90^\circ$) holds along the $z$-axis (diagonal states) and the equatorial plane (equal populations).  The angle drops below $90^\circ$ only in the interior regions where both coherence and population imbalance are nonzero.}
\end{figure}

\section*{Why the angle matters beyond geometry}

The angle has a concrete statistical meaning.

Suppose you want to estimate a classical property of the state---say, the population of one energy level---but the state also carries coherence that you cannot directly control or measure.  In statistical estimation, unknown quantities that affect your measurement but are not what you are trying to estimate are called \emph{nuisance parameters}.

The paper proves that the Bures angle exactly determines how much estimation precision you lose to the coherence nuisance.  Specifically, $\sin^2\theta_B$ is the fraction of \emph{quantum Fisher information} (the fundamental quantum limit on estimation precision, set by the Cram\'er--Rao bound) for the population parameter that survives after the coherence parameters are eliminated.

When $\theta_B = 90^\circ$ (perpendicular), you lose nothing: the coherence nuisance does not interfere with population estimation at all.  When $\theta_B < 90^\circ$ (oblique), some Fisher information is lost.  The formula tells you exactly how much.

This converts the geometric angle into something an experimentalist cares about: how much does quantum structure cost you when you are trying to learn something classical?

\section*{What changes in higher dimensions?}

Qubits ($d = 2$) are special.  In dimension two, equal populations can give perpendicularity even when coherence is present.

For $d \geq 3$ (qutrits and above), the paper shows the situation is much stricter.

\subsection*{The spectator argument}

Imagine three energy levels.  Levels 1 and 2 are coherent (their off-diagonal entry is nonzero), but level 3 is uninvolved---it has no coherence with either 1 or 2.  Level 3 is the ``spectator.''

The paper shows that the spectator's contribution to the metric stays exactly zero, while the coherent levels' contributions become nonzero.  That asymmetry---zero at the spectator, nonzero at the coherent levels---means the population and coherence directions can no longer be perpendicular.

This works in every dimension $d \geq 3$ and for any state where at least one coherence pair has an uninvolved third level.

\subsection*{Near the classical simplex}

Even without a spectator, close to diagonal states any nonzero coherence breaks perpendicularity.  The paper gives an explicit perturbative formula for how large the departure is.

\subsection*{The open question}

The remaining case is \emph{fully coherent} states, where every pair of levels is coherent and there is no uninvolved spectator.  The paper conjectures that perpendicularity still fails, but cannot prove it.

\begin{figure}[ht]
\centering
\includegraphics[width=0.92\textwidth]{figures/fig6_qudit_angles.pdf}
\caption{Numerical evidence for the higher-dimensional conjecture.  Blue histograms: minimum angle between population and coherence directions for 2,000 random states per dimension.  Solid orange lines: the best angle found by adversarial optimization (100 starts per dimension, actively trying to push the angle toward $90^\circ$).  The dashed red line marks $90^\circ$ (perfect perpendicularity).  Neither random sampling nor deliberate optimization finds a non-diagonal state achieving perpendicularity.}
\end{figure}

\section*{What about thermodynamics?}

The paper is careful not to overstate the thermodynamic implications.  What it shows is:
\begin{itemize}[leftmargin=1.5em]
\item dephasing changes the \emph{non-equilibrium free energy} of the state (free energy is the amount of energy available to do useful work);
\item when the system's energy levels align with the dephasing basis, the free energy destroyed is exactly $k_B T\, C_r$, where $T$ is temperature and $C_r$ measures the amount of coherence;
\item this is precise free-energy accounting, not a claim that the destroyed coherence could have been extracted as work in practice (that is a more subtle question involving what resources are available).
\end{itemize}

The paper's headline contribution is the geometric structure, not the thermodynamics.

\section*{What the paper does \emph{not} claim}

The paper does not:
\begin{itemize}[leftmargin=1.5em]
\item explain why the universe has the dimensional structure it does---why quantum state spaces are high-dimensional and observers access low-dimensional projections.  The paper characterizes the geometric mechanism of the projection (from a $(d^2-1)$-dimensional quantum manifold onto a $(d-1)$-dimensional classical simplex), but takes this dimensional gap as given.  Why nature has this architecture is metaphysical, not empirically resolvable;
\item derive the Born rule (the rule for computing quantum probabilities---it is assumed);
\item explain why one specific measurement outcome occurs rather than another (the paper characterizes the space of possible outcomes, not the selection of a particular one);
\item prove the global higher-dimensional orthogonality conjecture for fully-coherent states;
\item claim that all thermodynamic questions about coherence are settled by one formula.
\end{itemize}

That restraint is important.  The paper's strength is that it says something precise and defensible inside standard quantum mechanics, without requiring new physics or interpretive commitments.

\section*{Why this may be worth reading if you are not a specialist}

Even if you have never encountered information geometry, the paper is trying to do something conceptually simple:

\begin{quote}
Take a famous qualitative idea---decoherence creates the classical world out of quantum structure---and state exactly what geometric structure is lost, what survives, and how cleanly the two types of information separate.
\end{quote}

The answer turns out to be richer than ``decoherence destroys quantum information.''  The classical/quantum separation is not binary---it is an angle that varies across state space and sharpens dynamically as decoherence proceeds.  The paper gives the exact formula for that angle and proves it has operational meaning for statistical estimation.

If the paper is right, decoherence is not just ``noise killing interference.''  It is a mathematically structured collapse from a richer geometry to a simpler one, and the geometry of that collapse is precisely computable.

\section*{If you want to read the full paper efficiently}

Suggested reading order:
\begin{enumerate}[leftmargin=1.5em]
\item the abstract and introduction;
\item the Bures separation angle theorem and its proof (the main result);
\item the estimation corollary (the operational payoff);
\item the spectator orthogonality criterion (the higher-dimensional extension);
\item the discussion section;
\item then return to the earlier geometric setup (Section~2) if you want the full derivations.
\end{enumerate}

Repository guide:
\begin{itemize}[leftmargin=1.5em]
\item \texttt{decoherence\_dimensional\_collapse.pdf}: the full manuscript;
\item \texttt{reader\_guide.pdf}: this note;
\item \texttt{figures/}: publication figures used in both documents;
\item \texttt{code/generate\_figures.py}: reproducible code for all figures and numerical experiments.
\end{itemize}

\vfill

\begin{center}
\small
This guide is intentionally pedagogical.  For the precise theorem statements, hypotheses, and proofs, consult the main manuscript.
\end{center}

\end{document}
